\section{Local mass of terminal success}\label{h:local}

At a terminal state, the output law is independent of the unknown
signal. A small spherical ball has small uniform mass, and the same is
true of the set of signals close to a fixed output. Combining these
two facts initializes the bound that will be propagated backward.

Write \(\sigma=\sigma_d\), and let \(B(z,r)\) denote a closed
Euclidean ball in the ambient space of its center. Vector norms are
Euclidean, and \(\|A\|_{\op}\) is the Euclidean operator norm. Constants
denoted by \(C\) are positive and absolute and may increase between
displays. For the main proof set
\begin{equation}\label{h:local:parameters}
 n=d-1,\qquad a=\frac n2,\qquad
 q=\left\lfloor\frac n8\right\rfloor.
\end{equation}
We work in sufficiently large dimension that
\(q\ge2\), \(q\ge d/16\), and \(a\ge d/3\). These are absolute
dimension requirements. The real number \(a\) is a local-mass
exponent, whereas the integer \(q\) is a projection-moment order.

For \(R>0\), say that a measurable \(h:S^{d-1}\to[0,1]\) has local
mass at scale \(R\) if
\begin{equation}\label{h:local:bound}
 \int_{B(z,r)}h\,d\sigma\le(Rr)^a
 \qquad(z\in S^{d-1},\ r>0).
\end{equation}
The scale is an upper-bound parameter; it is not a support radius.

\begin{lemma}[Spherical balls and terminal outputs]\label{h:local:caps}
For \(d\ge2\), \(z\in S^{d-1}\), and \(r>0\),
\begin{equation}\label{h:local:cap}
 \sigma_d(B(z,r))\le r^{d-1}.
\end{equation}
Consequently, if \(Q\) is any probability law on \(S^{d-1}\)
independent of the signal and
\[
 h(s)=Q\{v:\arccos\ip{v}{s}\le\epsilon\},
\]
then \(h\) has local mass at scale \(R=\epsilon\):
\begin{equation}\label{h:local:terminal}
 \int_{B(z,r)}h\,d\sigma_d\le(\epsilon r)^{(d-1)/2}.
\end{equation}
\end{lemma}

\begin{proof}
For \(0<r\le1\), rotate \(z\) to the north pole and write
\(n=d-1\). A point of \(S^{d-1}\cap B(z,r)\) has last coordinate at
least \(1-r^2/2\ge1/2\). On the upper hemisphere the surface-area
element as a graph over the first \(n\) coordinates is the reciprocal
of that last coordinate, hence at most two on this cap. Its projection
is contained in the \(n\)-ball of radius \(r\). The full sphere has
area at least twice the volume of the unit \(n\)-ball, since both
hemispheres project onto that ball with area element at least one.
The ratio is at most \(r^n\). For \(r>1\), the total mass one gives
the same inequality.

If the angular distance from a fixed unit output \(v\) is at most
\(\epsilon\), its chordal distance is
\(2\sin(\theta/2)\le\theta\le\epsilon\). Thus the global successful
mass for \(v\) is at most \(\epsilon^n\). Integrating this inequality
against \(Q\) gives the same global bound for \(h\). Independently,
\(0\le h\le1\) and \eqref{h:local:cap} give local mass at most \(r^n\).
Therefore
\[
 \int_{B(z,r)}h\,d\sigma_d
 \le\min\{r^n,\epsilon^n\}\le(\epsilon r)^{n/2},
\]
where the last inequality follows by considering \(r\le\epsilon\)
and \(r\ge\epsilon\).
\end{proof}

The two bounds play different roles: \(r^{d-1}\) controls
concentration in a small region, and \(\epsilon^{d-1}\) controls the
total successful mass. Their geometric mean has the same exponent in
\(r\) and \(\epsilon\), which is the form preserved by one block.
